\documentclass[twocolumn,prd,nofootinbib]{revtex4-2}
\usepackage{amsmath,amssymb,graphicx,hyperref,booktabs,algorithm,algpseudocode}

\begin{document}

\title{NEXUS: A Domain-Agnostic Self-Learning Daemon for\\Automated Parameter Exploration and Anomaly Detection}

\author{Danny Lee Eldridge}
\affiliation{Independent Researcher, United States}
\email{dannyeldridge@proton.me}

\date{\today}

\begin{abstract}
We introduce \textsc{NEXUS} (Neural Exploration eXplorer with Unified Scanning), a domain-agnostic optimization daemon that autonomously explores parameter spaces, detects anomalies, and identifies discoveries without human intervention. The system combines adaptive grid scanning, Nelder-Mead simplex optimization, smart pathfinding with exploration-exploitation balancing, and four distinct anomaly detectors (constraint violations, statistical outliers, gradient spikes, and plateau detection). Unlike domain-specific tools, \textsc{NEXUS} operates on any model expressible as a fitness function over tunable parameters. We demonstrate its effectiveness on classical optimization benchmarks (Rosenbrock, Rastrigin), hyperparameter tuning, and engineering design problems. The system achieves near-optimal solutions on Rosenbrock ($f < 0.01$) within 2000 evaluations and correctly identifies multi-modal structure in Rastrigin. Zero runtime dependencies and a TypeScript implementation ensure deployment in under 10 seconds on any platform.
\end{abstract}

\maketitle

\section{Introduction}

Parameter optimization and exploration are fundamental to virtually every quantitative discipline: machine learning hyperparameter tuning \cite{bergstra2012}, engineering design optimization \cite{deb2002}, scientific model calibration \cite{kennedy2001}, and financial strategy development. While domain-specific tools exist (Optuna \cite{optuna}, Ray Tune \cite{raytune}, MATLAB's \texttt{fmincon}), no single framework provides:

\begin{enumerate}
    \item Autonomous daemon operation (set-and-forget)
    \item Multi-phase adaptive exploration
    \item Integrated anomaly detection
    \item Discovery logging (degeneracies, multi-modal landscapes)
    \item Smart pathfinding to avoid redundant exploration
    \item Zero dependencies and instant deployment
\end{enumerate}

\textsc{NEXUS} fills this gap as a lightweight, domain-agnostic daemon that continuously refines its understanding of any parameter space.

\section{Algorithm}

\subsection{Six-Phase Exploration}

\begin{algorithm}[h]
\caption{NEXUS Adaptive Scan}
\begin{algorithmic}[1]
\State \textbf{Phase 1: Grid Scan} — Coarse sampling (full grid or LHS)
\State \textbf{Phase 2: Adaptive Zoom} — $k$ passes narrowing by factor $f$
\For{$i = 1$ to $k$}
    \State Compute weighted centroid: $\bar{\theta} = \sum w_i \theta_i / \sum w_i$
    \State $w_i = \exp(-\frac{1}{2}(f_i - f_\mathrm{min}))$
    \State Narrow ranges: $[\bar{\theta} - f \cdot \sigma, \bar{\theta} + f \cdot \sigma]$
    \State Re-scan narrowed region
\EndFor
\State \textbf{Phase 3: Nelder-Mead} — Simplex optimization from best point
\State \textbf{Phase 4: Smart Explore} — Sample unvisited regions
\State \textbf{Phase 5: Anomaly Detection} — Four detectors (Sec.\ \ref{sec:anomaly})
\State \textbf{Phase 6: Discovery Logging} — Correlations, degeneracies
\end{algorithmic}
\end{algorithm}

\subsection{Smart Pathfinding}

The exploration module maintains a visited-region map by discretizing parameter space into $10^d$ bins (where $d$ is dimensionality). New candidate points are scored by \emph{novelty}:
\begin{equation}
    \mathrm{novelty}(\theta) = \frac{1}{1 + n_\mathrm{visits}(\mathrm{bin}(\theta))}
\end{equation}

An exploration rate $\epsilon \in [0, 1]$ controls the fraction of evaluations dedicated to novel regions versus refining known-good regions (exploitation). This implements a simple bandit-like explore-exploit balance.

\subsection{Anomaly Detection}\label{sec:anomaly}

Four independent detectors operate simultaneously:

\begin{enumerate}
    \item \textbf{Constraint violations:} Domain-specific rules (physics laws, engineering limits) checked at every evaluation point.
    \item \textbf{Statistical outliers:} Points with $|Z| > \tau$ (configurable threshold, default $\tau = 3$) in the fitness distribution.
    \item \textbf{Gradient spikes:} Adjacent points where $|\nabla f| > 10 \bar{f}$, indicating sharp discontinuities or phase transitions.
    \item \textbf{Plateau detection:} Regions where the top 10\% of results span $<0.1\%$ fitness variation but $>30\%$ parameter range, indicating unconstrained directions.
\end{enumerate}

\subsection{Anomaly Chasing}

High-severity anomalies ($s > 0.5$) trigger automatic fine-grained exploration: 10 additional evaluations within 2\% of the anomaly's parameter values. This reveals whether anomalies are isolated numerical artifacts or genuine features of the landscape.

\subsection{Discovery Engine}

Pearson correlation coefficients are computed for all parameter pairs. Pairs with $|r| > 0.85$ are flagged as \emph{degeneracies} — indicating the model is overparameterized in that direction. Multi-modal landscapes are identified when multiple well-separated minima ($d_\mathrm{norm} > 0.2$) achieve near-optimal fitness ($f < 1.1 f_\mathrm{best}$).

\section{Daemon Architecture}

The daemon wraps the scanner in an infinite loop with:
\begin{itemize}
    \item \textbf{Progressive resolution:} Grid density increases from $n_0$ to $n_\mathrm{max}$ across cycles
    \item \textbf{Cumulative tracking:} Global best-fit, total anomalies, discoveries persist across cycles
    \item \textbf{Automatic reporting:} JSON + Markdown reports generated each cycle
    \item \textbf{Fitness targets:} Optional early stopping when fitness $\leq f_\mathrm{target}$
    \item \textbf{Pro dashboard:} Real-time console output with Unicode visualizations
\end{itemize}

\section{Results}

\subsection{Rosenbrock Function}

The Rosenbrock function $f(x,y) = (1-x)^2 + 100(y-x^2)^2$ has a global minimum at $(1,1)$ inside a narrow curved valley. \textsc{NEXUS} achieves $f < 0.01$ within 2000 evaluations (5 zoom passes + Nelder-Mead), correctly identifying the valley structure.

\subsection{Rastrigin Function}

The Rastrigin function is highly multi-modal ($\sim 100$ local minima in $[-5.12, 5.12]^2$). \textsc{NEXUS} finds $f < 1$ (near the global minimum at origin) and correctly logs multi-modal landscape as a discovery.

\subsection{ML Hyperparameter Tuning}

On a simulated neural network loss landscape (learning rate $\times$ batch size $\times$ dropout), \textsc{NEXUS} identifies the optimal region and flags a constraint violation: high learning rate combined with low dropout produces unstable training.

\subsection{Engineering Design}

For cantilever beam weight minimization subject to stress ($< 250$ MPa) and deflection ($< L/200$) constraints, \textsc{NEXUS} finds feasible designs and logs all constraint-violating regions for engineer review.

\section{Comparison with Existing Tools}

\begin{table}[h]
\centering
\begin{tabular}{lccccc}
\toprule
Feature & NEXUS & Optuna & Ray & MATLAB \\
\midrule
Daemon mode & \checkmark & $\times$ & $\times$ & $\times$ \\
Anomaly detection & \checkmark & $\times$ & $\times$ & $\times$ \\
Smart pathfinding & \checkmark & $\sim$ & $\sim$ & $\times$ \\
Zero dependencies & \checkmark & $\times$ & $\times$ & $\times$ \\
Domain agnostic & \checkmark & ML only & ML only & \checkmark \\
Setup time & 10s & 5min & 30min & hours \\
\bottomrule
\end{tabular}
\caption{Feature comparison with existing optimization tools.}
\end{table}

\section{Software Availability}

\textsc{NEXUS} is implemented in TypeScript with zero runtime dependencies. It requires Node.js $\geq 18$ and installs via \texttt{npm install nexus-daemon}. The package includes 14 passing tests, two worked examples, and comprehensive documentation.

Contact: \texttt{dannyeldridge@proton.me}

\section{Conclusions}

\textsc{NEXUS} demonstrates that a simple, well-designed daemon architecture can provide autonomous parameter exploration with anomaly detection across any domain. The combination of adaptive scanning, smart pathfinding, and four anomaly detectors produces insights that go beyond simple optimization — revealing degeneracies, instabilities, and landscape structure that inform model design decisions.

\begin{thebibliography}{10}
\bibitem{bergstra2012} Bergstra, J.\ \& Bengio, Y.\ 2012, JMLR, 13, 281
\bibitem{deb2002} Deb, K.\ et al.\ 2002, IEEE TEC, 6, 182
\bibitem{kennedy2001} Kennedy, M.\ C.\ \& O'Hagan, A.\ 2001, JRSSB, 63, 425
\bibitem{optuna} Akiba, T.\ et al.\ 2019, KDD, 2623
\bibitem{raytune} Liaw, R.\ et al.\ 2018, arXiv:1807.05118
\bibitem{neldermead} Nelder, J.\ A.\ \& Mead, R.\ 1965, Computer J., 7, 308
\end{thebibliography}

\end{document}
